\chapter{Algebraic foundations}
\label{chap:algebraic-foundations}

This flat chapter is the first layer of the Blueprint.  It contains only the exact, graded, derived, and filtered algebra needed to state later spectral-sequence constructions; it contains no external mathematical input.

\section{Exact, graded, and derived algebra}

\begin{definition}[Graded object]
  \label{def:mathlib-graded-object}
  \lean{CategoryTheory.GradedObject}
  \mathlibok
  For an index type $I$ and a category $\mathcal C$, a graded object is a
  family $(A_i)_{i\in I}$ of objects of $\mathcal C$.  In the pinned Mathlib
  revision this is definitionally the function type $I\to\mathcal C$.
\end{definition}

% SOURCE: Mathlib/Data/ZMod/Defs.lean:140--144, declaration ZMod; Mathlib/Algebra/Field/ZMod.lean:29--30, prime-field instance.
\begin{definition}[The $\operatorname{ZMod}$ coefficient family]
  \label{def:mathlib-f2}
  \lean{ZMod}
  \mathlibok
  Mathlib's declaration is the family $\operatorname{ZMod}(n)$ indexed by
  $n\in\mathbb N$; it does not itself fix a modulus.  This node records only
  that generic family.  Later sign simplifications must use an explicit
  characteristic-two specialization rather than treating every abelian
  category as characteristic two.
\end{definition}

% SOURCE: Mathlib/Algebra/Category/ModuleCat/Basic.lean:48--62, declaration ModuleCat.
\begin{definition}[Category of modules over a ring]
  \label{def:mathlib-f2-module-category}
  \lean{ModuleCat}
  \mathlibok
  Mathlib's declaration $\operatorname{ModuleCat}(R)$ is the category of
  modules over an arbitrary ring $R$.  This generic declaration does not
  silently choose the project's coefficient ring.
\end{definition}

% VERIFIED-IN: KIP126/Core/Algebra/Coefficients.lean, declaration KIP126.Core.Algebra.F2.
\begin{definition}[Project mod--2 coefficient field]
  \label{def:core-f2}
  \lean{KIP126.Core.Algebra.F2}
  \leanok
  \uses{def:mathlib-f2}
  The project fixes the abbreviation
  $\mathbb F_2:=\operatorname{ZMod}(2)$.  Its pinned Mathlib instances provide
  $1+1=0$, characteristic two, decidable equality, and finiteness.
\end{definition}

% VERIFIED-IN: KIP126/Core/Algebra/Coefficients.lean, declaration KIP126.Core.Algebra.F2ModuleCat.
\begin{definition}[Project mod--2 coefficient category]
  \label{def:core-f2-module-category}
  \lean{KIP126.Core.Algebra.F2ModuleCat}
  \leanok
  \uses{def:core-f2,def:mathlib-f2-module-category}
  The coefficient category used by the compiled classical and synthetic Adams
  interfaces is the explicit abbreviation
  $\operatorname{ModuleCat}(\mathbb F_2)$.
\end{definition}

% SOURCE: Mathlib/Algebra/Category/ModuleCat/Abelian.lean:74--78, instance ModuleCat.abelian.
\begin{theorem}[Module categories are abelian]
  \label{thm:mathlib-module-category-abelian}
  \lean{ModuleCat.abelian}
  \mathlibok
  \uses{def:mathlib-f2-module-category}
  For every ring $R$, the category $\operatorname{ModuleCat}(R)$ carries
  Mathlib's abelian-category instance.  In particular this applies after the
  project specializes $R$ to $\mathbb F_2$.
\end{theorem}

% SOURCE: Mathlib/CategoryTheory/Abelian/Basic.lean, declaration CategoryTheory.Abelian.
\begin{definition}[Abelian category]
  \label{def:mathlib-abelian}
  \lean{CategoryTheory.Abelian}
  \mathlibok
  In the pinned Mathlib declaration, an abelian category is preadditive, has
  kernels, cokernels, and finite products, and every monomorphism and
  epimorphism is normal.  A zero object and the usual finite biproduct
  consequences are derived from these fields.
\end{definition}

% SOURCE: Mathlib/Algebra/Homology/HomologicalComplex.lean, declaration CategoryTheory.HomologicalComplex.
\begin{definition}[Homological complex]
  \label{def:mathlib-homological-complex}
  \lean{HomologicalComplex}
  \mathlibok
  For a category $\mathcal C$, an index type $\iota$, and a complex shape
  $c$, a homological complex consists of objects $K_i$, differentials
  $d_{ij}:K_i\to K_j$ allowed by $c$, and the square-zero relations required
  by that shape.
\end{definition}

% SOURCE: Mathlib/Algebra/Homology/HomologicalComplex.lean, abbreviation CategoryTheory.ChainComplex.
\begin{definition}[Chain complex]
  \label{def:mathlib-chain-complex}
  \lean{ChainComplex}
  \mathlibok
  \uses{def:mathlib-homological-complex}
  A chain complex in $\mathcal C$ indexed by an additive index type is the
  Mathlib homological complex with differential of degree $-1$.
\end{definition}

% SOURCE: Mathlib/Algebra/Homology/HomologicalComplex.lean, abbreviation CategoryTheory.CochainComplex.
\begin{definition}[Cochain complex]
  \label{def:mathlib-cochain-complex}
  \lean{CochainComplex}
  \mathlibok
  \uses{def:mathlib-homological-complex}
  A cochain complex in $\mathcal C$ indexed by an additive index type is the
  Mathlib homological complex with differential of degree $+1$.  Passing from a
  chain complex to a cochain complex therefore requires an explicit reindexing;
  the two abbreviations are not interchangeable.
\end{definition}

% SOURCE: Mathlib/Algebra/Homology/ShortComplex/Exact.lean, declaration CategoryTheory.ShortComplex.Exact.
\begin{definition}[Exact short complex]
  \label{def:mathlib-short-complex-exact}
  \lean{CategoryTheory.ShortComplex.Exact}
  \mathlibok
  For a short complex $X_1\xrightarrow{f}X_2\xrightarrow{g}X_3$, Mathlib
  exactness means that there exists homology data whose homology object is
  zero.  In an abelian category, the next node identifies this with the
  familiar image-to-kernel comparison being an isomorphism.
\end{definition}

% SOURCE: Mathlib/CategoryTheory/Abelian/Exact.lean:77, theorem CategoryTheory.ShortComplex.exact_iff_isIso_imageToKernel.
\begin{theorem}[Exactness via image and kernel]
  \label{thm:mathlib-exact-iff-image-kernel}
  \lean{CategoryTheory.ShortComplex.exact_iff_isIso_imageToKernel}
  \mathlibok
  \uses{def:mathlib-short-complex-exact,def:mathlib-abelian}
  In an abelian category, a short complex is exact if and only if its canonical
  morphism from the image of the first map to the kernel of the second map is
  an isomorphism.
\end{theorem}

% SOURCE: Mathlib/CategoryTheory/GradedObject.lean, declaration CategoryTheory.GradedObject.comap.
\begin{definition}[Regrading a Mathlib graded object]
  \label{def:mathlib-graded-comap}
  \lean{CategoryTheory.GradedObject.comap}
  \mathlibok
  \uses{def:mathlib-graded-object}
  A function $h:J\to I$ induces Mathlib's regrading functor from
  $I$-graded objects to $J$-graded objects by precomposition.  This declaration
  regrades objects only; it does not by itself regrade complexes or spectral
  sequences.
\end{definition}

% SOURCE: Mathlib/Algebra/Homology/ComplexShape.lean, declaration ComplexShape.up'.
\begin{definition}[Complex shape of a fixed additive degree]
  \label{def:mathlib-complex-shape-up}
  \lean{ComplexShape.up'}
  \mathlibok
  For an additive index type and a fixed degree $a$, Mathlib's
  $\operatorname{up}'(a)$ permits a differential from $i$ to $j$ exactly when
  $i+a=j$.  The AIM Adams shape therefore uses the integer pair
  $(r,r-1)$, not Mathlib's prepackaged cohomological convention $(r,1-r)$.
\end{definition}

% SOURCE: Mathlib/Algebra/Homology/DerivedCategory/Ext/Basic.lean, declaration CategoryTheory.Abelian.Ext.
\begin{definition}[Mathlib derived Ext]
  \label{def:mathlib-derived-ext}
  \lean{CategoryTheory.Abelian.Ext}
  \mathlibok
  \uses{def:mathlib-abelian}
  In an abelian category satisfying Mathlib's Ext-size hypothesis,
  $\operatorname{Ext}(X,Y,n)$ is a singly graded derived Ext type indexed by
  $n\in\mathbb N$.  It does not provide the Steenrod algebra, internal grading,
  comodules, or the bigraded Adams $\operatorname{Ext}_{\mathcal A}^{s,t}$.
\end{definition}

% SOURCE: Mathlib/CategoryTheory/Shift/Basic.lean:60--68, declaration CategoryTheory.HasShift.
\begin{definition}[Mathlib shift structure]
  \label{def:mathlib-shift}
  \lean{CategoryTheory.HasShift}
  \mathlibok
  Mathlib's declaration supplies a coherent family of endofunctors indexed by
  an arbitrary additive monoid $A$.  The generic monoid-indexed structure does
  not assert that every shift is an autoequivalence.  The project uses
  $A=\mathbb Z$, where inverse indices supply the equivalence behaviour needed
  by triangulated constructions.
\end{definition}

% SOURCE: Mathlib/CategoryTheory/Triangulated/Basic.lean, declaration CategoryTheory.Pretriangulated.Triangle.
\begin{definition}[Mathlib triangle]
  \label{def:mathlib-triangle}
  \lean{CategoryTheory.Pretriangulated.Triangle}
  \mathlibok
  \uses{def:mathlib-shift}
  A Mathlib triangle consists of $X\to Y\to Z\to X[1]$ in a category with
  integer shifts.
\end{definition}

% SOURCE: Mathlib/CategoryTheory/Triangulated/Pretriangulated.lean, declaration CategoryTheory.Pretriangulated.
\begin{definition}[Mathlib pretriangulated structure]
  \label{def:mathlib-pretriangulated}
  \lean{CategoryTheory.Pretriangulated}
  \mathlibok
  \uses{def:mathlib-triangle}
  Mathlib's pretriangulated class chooses distinguished triangles and their
  closure laws.  It does not include a stable-spectrum model or compatibility
  of a monoidal product with distinguished triangles.
\end{definition}

% SOURCE: Mathlib/CategoryTheory/Triangulated/HomologicalFunctor.lean:82, declaration CategoryTheory.Functor.IsHomological.
\begin{definition}[Mathlib homological functor]
  \label{def:mathlib-homological-functor}
  \lean{CategoryTheory.Functor.IsHomological}
  \mathlibok
  \uses{def:mathlib-pretriangulated,def:mathlib-abelian}
  A functor from a pretriangulated category to an abelian category is
  homological when it preserves zero morphisms and sends every distinguished
  triangle to an exact short complex.
\end{definition}

% SOURCE: Mathlib/CategoryTheory/Shift/ShiftSequence.lean:48, declaration CategoryTheory.Functor.ShiftSequence.
\begin{definition}[Mathlib shift sequence for a functor]
  \label{def:mathlib-homological-shift-sequence}
  \lean{CategoryTheory.Functor.ShiftSequence}
  \mathlibok
  \uses{def:mathlib-shift}
  For a functor $F:\mathcal C\to\mathcal A$, a shift sequence indexed by
  $\mathbb Z$ chooses functors $F_n:\mathcal C\to\mathcal A$ and coherent
  isomorphisms identifying $F_n\circ[m]$ with $F_{m+n}$.  This is additional
  data on a homological functor and is what lets Mathlib state its long exact
  sequence with separately indexed functors $F_n$.
\end{definition}

% SOURCE: Mathlib/CategoryTheory/Triangulated/HomologicalFunctor.lean:173, declaration CategoryTheory.Functor.homologySequenceδ.
\begin{definition}[Homology-sequence connecting morphism]
  \label{def:mathlib-homology-sequence-connecting}
  \lean{CategoryTheory.Functor.homologySequenceδ}
  \mathlibok
  \uses{def:mathlib-triangle,def:mathlib-homological-shift-sequence}
  Let $F:\mathcal C\to\mathcal A$ have a $\mathbb Z$-shift sequence.  For a
  triangle $X\to Y\to Z\to X[1]$ and $n_0+1=n_1$, Mathlib constructs the
  connecting morphism $F_{n_0}(Z)\to F_{n_1}(X)$ by applying the shift map of
  $F$ to the third morphism of the triangle.
\end{definition}

% SOURCE: Mathlib/CategoryTheory/Triangulated/HomologicalFunctor.lean:181, declaration CategoryTheory.Functor.homologySequenceδ_naturality.
\begin{theorem}[Naturality of the homology-sequence connecting morphism]
  \label{thm:mathlib-homology-sequence-connecting-naturality}
  \lean{CategoryTheory.Functor.homologySequenceδ_naturality}
  \mathlibok
  \uses{def:mathlib-homology-sequence-connecting}
  For every morphism of triangles, the third component mapped by $F_{n_0}$,
  the two connecting morphisms, and the first component mapped by $F_{n_1}$
  form a commutative square.
\end{theorem}

% SOURCE: Mathlib/CategoryTheory/Triangulated/HomologicalFunctor.lean:217, declaration CategoryTheory.Functor.homologySequence_exact₂.
\begin{theorem}[Middle exactness of a homology sequence]
  \label{thm:mathlib-homology-sequence-exact-two}
  \lean{CategoryTheory.Functor.homologySequence_exact₂}
  \mathlibok
  \uses{def:mathlib-homological-functor,def:mathlib-triangle,def:mathlib-homological-shift-sequence}
  A Mathlib homological functor sends the middle two maps of every shifted
  distinguished triangle to an exact short complex.
\end{theorem}

% SOURCE: Mathlib/CategoryTheory/Triangulated/HomologicalFunctor.lean:226, declaration CategoryTheory.Functor.homologySequence_exact₃.
\begin{theorem}[Connecting-map exactness of a homology sequence]
  \label{thm:mathlib-homology-sequence-exact-three}
  \lean{CategoryTheory.Functor.homologySequence_exact₃}
  \mathlibok
  \uses{def:mathlib-homological-functor,def:mathlib-triangle,def:mathlib-homological-shift-sequence,def:mathlib-homology-sequence-connecting}
  Rotation gives exactness at the term adjacent to the homology-sequence
  connecting morphism.
\end{theorem}

% SOURCE: Mathlib/CategoryTheory/Triangulated/HomologicalFunctor.lean:234, declaration CategoryTheory.Functor.homologySequence_exact₁.
\begin{theorem}[Shifted exactness of a homology sequence]
  \label{thm:mathlib-homology-sequence-exact-one}
  \lean{CategoryTheory.Functor.homologySequence_exact₁}
  \mathlibok
  \uses{def:mathlib-homological-functor,def:mathlib-triangle,def:mathlib-homological-shift-sequence,def:mathlib-homology-sequence-connecting}
  Inverse rotation gives the remaining adjacent exactness statement, allowing
  the three short complexes to splice into a long exact sequence.
\end{theorem}

% SOURCE: Mathlib/Algebra/Homology/HomotopyCategory/HomologicalFunctor.lean:34, instance HomotopyCategory.instIsHomologicalIntUpHomologyFunctor.
\begin{theorem}[Cochain homology is homological on the homotopy category]
  \label{thm:mathlib-homotopy-cochain-homology-is-homological}
  \lean{HomotopyCategory.instIsHomologicalIntUpHomologyFunctor}
  \mathlibok
  \uses{def:mathlib-cochain-complex,def:mathlib-homological-functor,def:mathlib-abelian}
  If $\mathcal C$ is abelian, then for every $n\in\mathbb Z$ the degree-$n$
  homology functor from the homotopy category of $\mathbb Z$-indexed cochain
  complexes in $\mathcal C$ to $\mathcal C$ is homological.
\end{theorem}

% SOURCE: Mathlib/Algebra/Homology/HomotopyCategory/ShiftSequence.lean:192-196, instance HomotopyCategory.instShiftSequenceIntUpHomologyFunctorOfNat.
\begin{definition}[Shift sequence on cochain homology]
  \label{def:mathlib-homotopy-cochain-homology-shift-sequence}
  \lean{HomotopyCategory.instShiftSequenceIntUpHomologyFunctorOfNat}
  \mathlibok
  \uses{def:mathlib-cochain-complex,def:mathlib-homological-shift-sequence}
  In a preadditive category equipped with Mathlib's homology structure, the
  degree-zero homology functor on the homotopy category of $\mathbb Z$-indexed
  cochain complexes has a $\mathbb Z$-shift sequence whose degree-$n$ member is
  the degree-$n$ homology functor.
\end{definition}

% SOURCE: Mathlib/CategoryTheory/Monoidal/Category.lean, declaration CategoryTheory.MonoidalCategory.
\begin{definition}[Mathlib monoidal category]
  \label{def:mathlib-monoidal-category}
  \lean{CategoryTheory.MonoidalCategory}
  \mathlibok
  Mathlib's monoidal category provides tensor product, unit, associator, and
  unitors with their coherence laws.
\end{definition}

% SOURCE: Mathlib/CategoryTheory/Monoidal/Braided/Basic.lean, declaration CategoryTheory.SymmetricCategory.
\begin{definition}[Mathlib symmetric category]
  \label{def:mathlib-symmetric-category}
  \lean{CategoryTheory.SymmetricCategory}
  \mathlibok
  \uses{def:mathlib-monoidal-category}
  A symmetric category extends a monoidal category with a symmetric braiding.
  Exactness of tensoring a triangle remains extra project data.
\end{definition}

\section{Filtered objects and complexes}

\begin{definition}[Decreasing filtration]
  \leanok
  \label{def:core-filtration}
  \lean{KIP126.Core.Algebra.Filtration}
  \uses{def:mathlib-graded-object}
  Let $A=(A_i)_{i\in I}$ be a graded object in a category.  A decreasing
  filtration $F$ of $A$ assigns a subobject $F^sA_i\subseteq A_i$ for every
  $s\in\mathbb Z$ and $i\in I$, together with inclusions
  \[
    F^{s+1}A_i\subseteq F^sA_i.
  \]
\end{definition}

\begin{definition}[Arbitrary-level filtration inclusions]
  \leanok
  \label{def:core-filtration-inclusions}
  \lean{KIP126.Core.Algebra.Filtration.le_of_le}
  \lean{KIP126.Core.Algebra.Filtration.inclusion}
  \lean{KIP126.Core.Algebra.Filtration.inclusion_arrow}
  \lean{KIP126.Core.Algebra.Filtration.inclusion_refl}
  \lean{KIP126.Core.Algebra.Filtration.inclusion_comp}
  \lean{KIP126.Core.Algebra.Filtration.transport_arrow}
  \lean{KIP126.Core.Algebra.Filtration.inclusion_naturality}
  \uses{def:core-filtration}
  The successor inclusions of a decreasing filtration extend canonically to
  every pair of levels: for $t\le s$ there is a map
  $F^sA_i\to F^tA_i$.  These maps are identity at equal levels and compose
  along chains of inequalities.  The transport and naturality equations keep
  the construction well-typed when a graded index is represented by
  propositionally equal integer expressions.
\end{definition}

\begin{definition}[Associated graded piece]
  \leanok
  \label{def:core-associated-graded}
  \lean{KIP126.Core.Algebra.Filtration.associatedGraded}
  \uses{def:core-filtration,def:mathlib-abelian}
  If $\mathcal C$ is abelian and $F$ is a decreasing filtration of $A$, define
  \[
    \operatorname{gr}^s_F(A_i) := F^sA_i/F^{s+1}A_i
  \]
  as the cokernel of the canonical inclusion
  $F^{s+1}A_i\hookrightarrow F^sA_i$.
\end{definition}

\begin{definition}[Quotient to an associated graded piece]
  \leanok
  \label{def:core-to-associated-graded}
  \lean{KIP126.Core.Algebra.Filtration.toAssociatedGraded}
  \uses{def:core-associated-graded}
  For every $(s,i)$, the canonical quotient morphism is
  \[
    F^sA_i\longrightarrow \operatorname{gr}^s_F(A_i).
  \]
\end{definition}

\begin{definition}[Associated graded object]
  \leanok
  \label{def:core-associated-graded-object}
  \lean{KIP126.Core.Algebra.Filtration.associatedGradedObject}
  \uses{def:core-associated-graded}
  The associated graded object is the $(\mathbb Z\times I)$-graded object
  whose $(s,i)$-component is $\operatorname{gr}^s_F(A_i)$.
\end{definition}

\begin{definition}[Degreewise eventually-top filtration]
  \leanok
  \label{def:core-exhaustive-filtration}
  \lean{KIP126.Core.Algebra.Filtration.IsExhaustive}
  \uses{def:core-filtration}
  The project predicate named \texttt{IsExhaustive} requires that for every
  $i$ there is an $s$ with $F^sA_i=A_i$.  Since the filtration is decreasing,
  all lower levels are then also the whole object.  This degreewise
  eventual-top condition is stronger than ordinary exhaustiveness stated only
  as $\bigcup_sF^sA_i=A_i$ or by a colimit.
\end{definition}

\begin{definition}[Degreewise eventually-bottom filtration]
  \leanok
  \label{def:core-eventually-zero-filtration}
  \lean{KIP126.Core.Algebra.Filtration.IsEventuallyZero}
  \uses{def:core-filtration,def:mathlib-abelian}
  The predicate \texttt{IsEventuallyZero} requires that for every $i$ there
  is an $s$ with $F^sA_i=0$; all higher levels then vanish as well.  This
  degreewise eventual-bottom condition is stronger than separatedness
  $\bigcap_sF^sA_i=0$.
\end{definition}

\begin{definition}[Bounded-below filtration]
  \leanok
  \label{def:core-bounded-below-filtration}
  \lean{KIP126.Core.Algebra.Filtration.IsBoundedBelow}
  \uses{def:core-filtration}
  A filtration is bounded below if each $i$ has an integer $l(i)$ such
  that $F^sA_i=A_i$ for $s\le l(i)$.
\end{definition}

\begin{definition}[Bounded-above filtration]
  \leanok
  \label{def:core-bounded-above-filtration}
  \lean{KIP126.Core.Algebra.Filtration.IsBoundedAbove}
  \uses{def:core-filtration,def:mathlib-abelian}
  In an abelian category, a filtration is bounded above if each $i$ has an
  integer $u(i)$ such that $F^sA_i=0$ for $u(i)\le s$.
\end{definition}

\begin{definition}[Degreewise bounded filtration]
  \leanok
  \label{def:core-bounded-filtration}
  \lean{KIP126.Core.Algebra.Filtration.IsBounded}
  \uses{def:core-bounded-below-filtration,def:core-bounded-above-filtration}
  In an abelian category, a filtration is degreewise bounded when it is both
  bounded below and bounded above, with chosen bounds satisfying
  $l(i)\le u(i)$ for every $i$.
\end{definition}

\begin{lemma}[A lower bound gives degreewise eventual-top]
  \leanok
  \label{lem:core-bounded-below-exhaustive}
  \lean{KIP126.Core.Algebra.Filtration.IsBoundedBelow.isExhaustive}
  \uses{def:core-bounded-below-filtration,def:core-exhaustive-filtration}
  A bounded-below filtration satisfies the project's strong
  \texttt{IsExhaustive} predicate, hence is degreewise eventually top.
\end{lemma}

\begin{proof}
  \leanok
  For a component $i$, use $l(i)$ as the eventual-top witness.  The
  defining endpoint equation gives the conclusion.
\end{proof}

\begin{lemma}[An upper bound gives eventual vanishing]
  \leanok
  \label{lem:core-bounded-above-eventually-zero}
  \lean{KIP126.Core.Algebra.Filtration.IsBoundedAbove.isEventuallyZero}
  \uses{def:core-bounded-above-filtration,def:core-eventually-zero-filtration}
  In an abelian category, a bounded-above filtration is eventually zero.
\end{lemma}

\begin{proof}
  \leanok
  For a component $i$, use $u(i)$ as the witness for eventual vanishing.  The
  defining endpoint equation gives the conclusion.
\end{proof}

\begin{definition}[Filtered morphism]
  \leanok
  \label{def:core-filtered-morphism}
  \lean{KIP126.Core.Algebra.FilteredMorphism}
  \uses{def:core-filtration}
  For filtered graded objects $(A,F)$ and $(B,G)$, a filtered morphism is a
  degree-preserving map $f:A\to B$ such that, for every $(s,i)$, its component
  $f_i$ restricts to a map $F^sA_i\to G^sB_i$.  Equivalently, the square of
  the two subobject inclusions and $f_i$ commutes.
\end{definition}

\begin{definition}[Map on associated graded pieces]
  \leanok
  \label{def:core-associated-graded-map}
  \lean{KIP126.Core.Algebra.FilteredMorphism.associatedGradedMap}
  \uses{def:core-filtered-morphism,def:core-associated-graded}
  A filtered morphism $f:(A,F)\to(B,G)$ induces, for each $(s,i)$, a morphism
  \[
    \operatorname{gr}^s_F(A_i)\longrightarrow
    \operatorname{gr}^s_G(B_i),
  \]
  obtained from the commuting square on the two consecutive filtration levels
  by the universal property of cokernels.
\end{definition}

\begin{lemma}[Projection equation for an associated-graded map]
  \leanok
  \label{lem:core-associated-graded-map-projection}
  \lean{KIP126.Core.Algebra.FilteredMorphism.toAssociatedGraded_comp_associatedGradedMap}
  \uses{def:core-filtered-morphism,def:core-associated-graded-map}
  The quotient projection from a filtration level to its associated-graded
  piece followed by the induced map agrees with the chosen filtration-level
  factorisation followed by the target quotient projection.
\end{lemma}

\begin{proof}
  \leanok
  This is the cokernel universal-property equation defining the induced map.
\end{proof}

\begin{definition}[Functorial filtered morphisms]
  \leanok
  \label{def:core-filtered-morphism-functoriality}
  \lean{KIP126.Core.Algebra.FilteredMorphism.id}
  \lean{KIP126.Core.Algebra.FilteredMorphism.comp}
  \lean{KIP126.Core.Algebra.FilteredMorphism.associatedGradedMap_id}
  \lean{KIP126.Core.Algebra.FilteredMorphism.associatedGradedMap_comp}
  \uses{def:core-filtered-morphism,def:core-associated-graded-map}
  Filtered morphisms have the identity and composition operations needed for
  categorical transport.  The induced maps on associated graded pieces
  preserve those operations, so the quotient construction is functorial in
  filtered maps.
\end{definition}

\begin{definition}[Filtered chain complex]
  \leanok
  \label{def:core-filtered-chain-complex}
  \lean{KIP126.Core.SpectralSequence.FilteredComplex}
  \uses{def:core-filtration,def:mathlib-chain-complex,def:mathlib-abelian}
  In an abelian category, a filtered chain complex is a Mathlib chain complex
  $(C_k,d_k)$ together with a decreasing filtration $F^sC_k$ of its underlying
  $\mathbb Z$-graded object such that every differential restricts to a map
  \[
    F^sC_k\longrightarrow F^sC_{k-1}.
  \]
\end{definition}

\begin{definition}[Associated graded differential]
  \leanok
  \label{def:core-associated-graded-differential}
  \lean{KIP126.Core.SpectralSequence.FilteredComplex.associatedGradedDifferential}
  \uses{def:core-filtered-chain-complex,def:core-associated-graded}
  The restriction of $d_k$ in a filtered chain complex descends to a morphism
  \[
    d^{\operatorname{gr}}_{s,k}:
    \operatorname{gr}^s_F(C_k)\longrightarrow
    \operatorname{gr}^s_F(C_{k-1}).
  \]
\end{definition}

\begin{theorem}[The associated graded differential squares to zero]
  \leanok
  \label{thm:core-associated-graded-differential-sq}
  \lean{KIP126.Core.SpectralSequence.FilteredComplex.associatedGradedDifferential_sq}
  \uses{def:core-associated-graded-differential}
  For every filtration degree $s$ and chain degree $k$,
  \[
    d^{\operatorname{gr}}_{s,k-1}\circ
    d^{\operatorname{gr}}_{s,k}=0.
  \]
\end{theorem}

\begin{proof}
  \leanok
  Precompose with the quotient epimorphism $F^sC_k\longrightarrow
  \operatorname{gr}^s_F(C_k)$.  The resulting composite factors through
  $d_{k-1}\circ d_k$, which is zero because $C$ is a chain complex.  The
  epimorphism then gives the asserted equality.
\end{proof}
